\chapter{The Registers}\label{cha:registers}\markboth{The Registers}{The Registers}

Everything the machine's devices do, they do through the I/O page. This
appendix is the register-level reference, and it is not written by hand:
\pth{doc/guide/mkregs.py} lifts it out of the machine's own headers at
build time, so what you read here is what the emulator was compiled
against. When a device gains a register, this appendix gains it in the
next build --- which is the only way a reference of this kind stays true.

Nothing in it is reworded: every word below is a word from a header.
What the generator decides is only which column a word belongs in --- the
address on the left, the register's name after it, the comment's own
sentence following. Where a comment's own spacing was doing the work of a
table, that table is kept exactly as it was written, in monospace, because
there it is the spacing that carries the meaning.

The voice is still the voice of working code: terse, occasionally
opinionated, and using the machine's own shorthand --- \verb|LE| for
little-endian, \verb|R:| and \verb|W:| for a register that reads and
writes differently, \verb|$| for hex. Addresses are given as an offset
inside the device's block where the block's base is obvious, and in full
where it is not.

Chapter~\ref{cha:io} is the map of which device lives where; this is what
is inside each one.

\input{generated/registers}
